% ===== PRÉAMBULE CANONIQUE — à coller VERBATIM en tête de CHAQUE .tex =====
\documentclass[11pt,a4paper]{article}
\usepackage[utf8]{inputenc}\usepackage[T1]{fontenc}\usepackage{lmodern}
\usepackage[french]{babel}
\usepackage{amsmath,amssymb,mathtools}\usepackage{icomma}
\usepackage{siunitx}\sisetup{locale=FR}  % \qty{}{}, \si{}, \ang{}
\usepackage{xcolor}\usepackage{colortbl}
\usepackage{tikz}\usepackage{pgfplots}\pgfplotsset{compat=1.18}
\usepackage[siunitx,europeanresistors]{circuitikz} % circuits (élec.)
\usepackage[most]{tcolorbox}\usepackage{enumitem}\usepackage{array,booktabs}
\usepackage[a4paper,margin=2.2cm]{geometry}
\usetikzlibrary{arrows.meta,calc,angles,quotes,positioning,patterns,%
  backgrounds,shapes.geometric,decorations.pathmorphing,decorations.markings,intersections}
\definecolor{primary}{RGB}{222,49,99}\definecolor{primarydark}{RGB}{150,22,55}
\definecolor{defcol}{RGB}{0,114,178}\definecolor{methcol}{RGB}{52,92,148}
\definecolor{excol}{RGB}{0,135,90}\definecolor{attncol}{RGB}{200,40,55}
\tcbset{boxbase/.style={enhanced,breakable,boxrule=0.9pt,arc=2.5pt,
  left=3mm,right=3mm,top=2.3mm,bottom=2.3mm,fonttitle=\bfseries,coltitle=white}}
% --- boîtes TUTORIEL (noms EXACTS, ne pas renommer) ---
\newtcolorbox{cle}[1][]{boxbase,colback=defcol!6,colframe=defcol,
  title={\textsc{À retenir}\ifx\relax#1\relax\relax\else\textmd{ : \itshape #1}\fi}}
\newtcolorbox{methode}[1][]{boxbase,colback=methcol!7,colframe=methcol,sharp corners,
  title={\textsc{Méthode}\ifx\relax#1\relax\relax\else\textmd{ --- \itshape #1}\fi}}
\newtcolorbox{exemplebox}[1][]{boxbase,colback=excol!5,colframe=excol,
  title={\textsc{Exemple}\ifx\relax#1\relax\relax\else\textmd{ : \itshape #1}\fi}}
\newtcolorbox{piege}[1][]{boxbase,colback=attncol!6,colframe=attncol,
  title={\textsc{Piège}\ifx\relax#1\relax\relax\else\textmd{ : \itshape #1}\fi}}
\newtcolorbox{remarque}[1][]{boxbase,colback=black!4,colframe=black!45,
  title={\textsc{Remarque}\ifx\relax#1\relax\relax\else\textmd{ : \itshape #1}\fi}}
% --- boîtes EXERCICE / CORRIGÉ (noms EXACTS, auto-comptées) ---
\newtcolorbox[auto counter]{exo}[1][]{boxbase,colback=primary!4,colframe=primary,
  title={\textsc{Exercice}~\thetcbcounter\ifx\relax#1\relax\relax\else\textmd{ --- \itshape #1}\fi}}
\newtcolorbox[auto counter]{corrige}[1][]{boxbase,colback=excol!4,colframe=excol,
  title={\textsc{Corrigé}~\thetcbcounter\ifx\relax#1\relax\relax\else\textmd{ --- \itshape #1}\fi}}
\newcommand{\vv}[1]{\vec{#1}}\newcommand{\ds}{\displaystyle}
\newcommand{\R}{\mathbb{R}}\newcommand{\N}{\mathbb{N}}\newcommand{\Z}{\mathbb{Z}}
% =========================================================================
\begin{document}

\begin{corrige}[hiérarchie des couleurs dans le verre]
\begin{enumerate}[label=\alph*)]
  \item $v=c/n$ : plus $n$ est grand, plus $v$ est petite. Comme $n_{\text{bleu}}>n_{\text{vert}}>n_{\text{rouge}}$ :
        \[ v_{\text{bleu}}<v_{\text{vert}}<v_{\text{rouge}}. \]
  \item La couleur la plus déviée est celle de plus grand indice : le \textbf{bleu} (il se rapproche
        le plus de la normale, $\hat r_{\text{bleu}}$ le plus petit).
\end{enumerate}
\end{corrige}

\begin{corrige}[une lumière rouge sur un prisme]
\begin{enumerate}[label=\alph*)]
  \item \textbf{Oui}, elle est déviée : tout rayon qui change de milieu est réfracté (dévié).
  \item \textbf{Non}, elle n'est pas décomposée : la dispersion sépare des \emph{couleurs différentes}.
        Une lumière monochromatique ne contient qu'une seule longueur d'onde — il n'y a rien à séparer.
\end{enumerate}
\end{corrige}

\begin{corrige}[la profondeur apparente d'une pièce]
\[ h'=\frac{n_2}{n_1}\,h=\frac{1}{1{,}33}\times\qty{1.60}{m}\approx\qty{1.20}{m}. \]
La pièce paraît à $\qty{1.20}{m}$ alors qu'elle est réellement à $\qty{1.60}{m}$ : elle semble
remontée d'environ $\qty{40}{cm}$.
\end{corrige}

\begin{corrige}[remonter à la profondeur réelle]
On inverse la relation $h'=\dfrac{n_2}{n_1}h$ :
\[ h=\frac{n_1}{n_2}\,h'=\frac{1{,}33}{1}\times\qty{0.60}{m}\approx\qty{0.80}{m}. \]
Le poisson est en réalité \emph{plus profond} ($\qty{0.80}{m}$) qu'il ne paraît : c'est pourquoi il
faut viser \emph{sous} le poisson que l'on voit.
\end{corrige}

\begin{corrige}[la latte brisée]
\begin{enumerate}[label=\alph*)]
  \item Elle paraît \textbf{plus proche} de la surface : les rayons issus de $A$ s'éloignent de la
        normale en sortant de l'eau, et leurs prolongements se croisent en $A'$, plus haut que $A$.
  \item Un objet immergé paraît globalement plus \textbf{gros} (et plus proche) qu'en réalité : ce
        n'est qu'une illusion d'optique due à la réfraction.
\end{enumerate}
\end{corrige}

\end{document}
